\DeclareEditionOverlay{CM-C02-U001}{%
\begin{EditionUnitCard}{Tres miradas de una sucesión}
Relaciona índice, término y regla antes de calcular.
\begin{EditionCheckpointBox}{Comprueba}
Describe una misma sucesión con palabras, varios términos y una regla cuando sea posible.
\end{EditionCheckpointBox}
\end{EditionUnitCard}}

\DeclareEditionOverlay{CM-C02-U002}{%
\begin{EditionUnitCard}{Del corredor a los cuantificadores}
Lee primero el corredor alrededor del límite y después la frase simbólica.
\RegisteredBookFigure{FIG-C02-001}{fig:c02:sequence-corridor}{Corredor epsilon--N para una sucesión convergente.}{Términos de una sucesión que, después de un índice N, permanecen dentro del corredor entre L menos epsilon y L más epsilon.}{\SequenceCorridorVisual}
\begin{EditionWarningBox}{Atención}
N puede depender de epsilon; no puede depender del índice n que viene después.
\end{EditionWarningBox}
\end{EditionUnitCard}}

\DeclareEditionOverlay{CM-C02-U003}{%
\begin{EditionUnitCard}{Primero decide qué teorema cabe}
Unicidad y acotación describen la convergencia; las reglas permiten combinar límites ya controlados.
\begin{EditionCheckpointBox}{Comprueba}
Antes de un cociente, identifica explícitamente el límite del denominador.
\end{EditionCheckpointBox}
\end{EditionUnitCard}}

\DeclareEditionOverlay{CM-C02-U004}{%
\begin{EditionUnitCard}{Una conexión de existencia}
La monotonía ordena los términos; la acotación impide escapar; la completitud cierra el argumento.
\begin{EditionCheckpointBox}{Comprueba}
Señala por separado dónde se prueba monotonía y dónde se obtiene una cota.
\end{EditionCheckpointBox}
\end{EditionUnitCard}}

\DeclareEditionOverlay{CM-C02-U005}{%
\begin{EditionUnitCard}{Mirar términos tardíos}
El criterio de Cauchy evita nombrar de antemano el posible límite.
\begin{EditionWarningBox}{Atención}
La proximidad se exige cuando ambos índices son suficientemente grandes.
\end{EditionWarningBox}
\end{EditionUnitCard}}

\DeclareEditionOverlay{CM-C02-U006}{%
\begin{EditionUnitCard}{Crecer sin cota no es converger en R}
Traduce cada símbolo de infinito a una afirmación con umbrales.
\begin{EditionCheckpointBox}{Comprueba}
Explica qué debe ocurrir después de un índice para superar cualquier cota dada.
\end{EditionCheckpointBox}
\end{EditionUnitCard}}

\DeclareEditionOverlay{CM-C02-U007}{%
\begin{EditionUnitCard}{Ampliación: criterio especializado}
Estúdialo solo después de dominar las reglas y estrategias nucleares.

\end{EditionUnitCard}}

\DeclareEditionOverlay{CM-C02-U008}{%
\begin{EditionUnitCard}{Aproximarse no es llegar}
Se observa f(x) cuando x se acerca a a por puntos admitidos del dominio.
\RegisteredBookFigure{FIG-C02-002}{fig:c02:epsilon-delta}{Lectura geométrica de epsilon--delta.}{Una franja vertical alrededor de a se envía dentro de una franja horizontal alrededor de L, excluyendo el punto central de aproximación.}{\EpsilonDeltaVisual}
\begin{EditionWarningBox}{Atención}
El límite puede existir aunque f(a) no exista o tenga otro valor.
\end{EditionWarningBox}
\end{EditionUnitCard}}

\DeclareEditionOverlay{CM-C02-U009}{%
\begin{EditionUnitCard}{Un puente entre dos lenguajes}
La caracterización secuencial prueba existencia solo cuando controla todas las sucesiones admisibles.
\begin{EditionCheckpointBox}{Comprueba}
Para refutar un límite, busca dos sucesiones con imágenes incompatibles.
\end{EditionCheckpointBox}
\end{EditionUnitCard}}

\DeclareEditionOverlay{CM-C02-U010}{%
\begin{EditionUnitCard}{El dominio decide por dónde acercarse}
Dibuja primero los puntos del dominio cercanos al punto de interés.
\begin{EditionCheckpointBox}{Comprueba}
Indica qué límites laterales tienen sentido antes de comparar sus valores.
\end{EditionCheckpointBox}
\end{EditionUnitCard}}

\DeclareEditionOverlay{CM-C02-U011}{%
\begin{EditionUnitCard}{Método antes que manipulación}
Clasifica la expresión, decide el resultado pertinente y conserva sus condiciones.
\begin{EditionCheckpointBox}{Comprueba}
Explica por qué el resultado elegido se aplica antes de hacer álgebra.
\end{EditionCheckpointBox}
\end{EditionUnitCard}}

\DeclareEditionOverlay{CM-C02-U012}{%
\begin{EditionUnitCard}{Una forma indeterminada no responde}
La etiqueta indica que hace falta transformar o estimar la expresión.
\begin{EditionWarningBox}{Atención}
No escribas el valor del límite a partir de la forma inicial.
\end{EditionWarningBox}
\end{EditionUnitCard}}

\DeclareEditionOverlay{CM-C02-U013}{%
\begin{EditionUnitCard}{Comparar cerca de un punto}
Toda equivalencia asintótica tiene un lugar y un alcance.
\begin{EditionCheckpointBox}{Comprueba}
Anota junto a cada equivalencia el punto al que se refiere.
\end{EditionCheckpointBox}
\end{EditionUnitCard}}

\DeclareEditionOverlay{CM-C02-U014}{%
\begin{EditionUnitCard}{Continuidad reúne tres piezas}
Pregunta por dominio, valor de la función y límite relativo al dominio.
\begin{EditionCheckpointBox}{Comprueba}
Comprueba cada pieza antes de declarar continuidad.
\end{EditionCheckpointBox}
\end{EditionUnitCard}}

\DeclareEditionOverlay{CM-C02-U015}{%
\begin{EditionUnitCard}{Las familias elementales conservan su dominio}
La continuidad no elimina puntos donde la expresión no está definida.
\begin{EditionWarningBox}{Atención}
Escribe el dominio antes de invocar continuidad elemental.
\end{EditionWarningBox}
\end{EditionUnitCard}}

\DeclareEditionOverlay{CM-C02-U016}{%
\begin{EditionUnitCard}{Diagnóstico antes que etiqueta}
Calcula límites laterales y revisa el valor antes de clasificar.
\begin{EditionCheckpointBox}{Comprueba}
Haz una tabla con límite izquierdo, derecho y valor de la función.
\end{EditionCheckpointBox}
\end{EditionUnitCard}}

\DeclareEditionOverlay{CM-C02-U017}{%
\begin{EditionUnitCard}{Un teorema de existencia}
El resultado garantiza extremos; no dice todavía dónde están.
\begin{EditionWarningBox}{Atención}
Revisa continuidad e intervalo cerrado y acotado antes de aplicarlo.
\end{EditionWarningBox}
\end{EditionUnitCard}}

\DeclareEditionOverlay{CM-C02-U018}{%
\begin{EditionUnitCard}{Existencia y aproximación}
Bolzano garantiza un cero; la bisección ayuda a localizarlo.
\begin{EditionCheckpointBox}{Comprueba}
Señala qué dato garantiza existencia y cuál reduce el intervalo.
\end{EditionCheckpointBox}
\end{EditionUnitCard}}

\DeclareEditionOverlay{CM-C02-U019}{%
\begin{EditionUnitCard}{Entre dos valores}
Comprueba qué valor objetivo queda entre las imágenes de los extremos.
\begin{EditionCheckpointBox}{Comprueba}
Explica por qué el teorema garantiza existencia pero no unicidad.
\end{EditionCheckpointBox}
\end{EditionUnitCard}}

